
\documentclass[11pt]{article}
\usepackage[a4paper,margin=0.9in]{geometry}
\usepackage[T1]{fontenc}
\usepackage[utf8]{inputenc}
\usepackage{lmodern}
\usepackage{amsmath,amssymb,amsthm,mathtools,mathrsfs}
\usepackage{microtype}
\usepackage{fancyhdr}
\usepackage{array,longtable,enumitem,multicol,booktabs}
\usepackage{tikz}
\setlength{\parindent}{1.5em}
\setlength{\parskip}{0.18em}
\setlength{\headheight}{14pt}
\emergencystretch=8em
\pagestyle{fancy}
\fancyhf{}
\cfoot{\thepage}
\providecommand{\mat}[4]{\begin{pmatrix}#1&#2\\#3&#4\end{pmatrix}}
\providecommand{\mmat}[9]{\begin{pmatrix}#1&#2&#3\\#4&#5&#6\\#7&#8&#9\end{pmatrix}}
\providecommand{\tg}{\operatorname{tg}}
\providecommand{\idxitem}[1]{\par\hangindent=1.4em\hangafter=1 #1}
\providecommand{\dd}{\,d}
\providecommand{\eps}{\varepsilon}
\providecommand{\abs}[1]{\left|#1\right|}
\providecommand{\Ftwo}{\mathbf F_2}
\providecommand{\Mod}{\operatorname{mod}}
\providecommand{\Det}[1]{\left|\begin{array}{#1}}
\providecommand{\E}{\mathrm E}
\providecommand{\D}{\Delta}
\providecommand{\Hs}{\mathcal H}
\providecommand{\Cov}{\operatorname{Cov}}

\usepackage[ngerman]{babel}
\lhead{Heinrich Weber, Lehrbuch der Algebra, Vol. II}
\rhead{German Source}
\begin{document}


\begin{center}
{\large Elfter Abschnitt.}\\[0.4em]
{\large Die Wendepunkte einer Curve dritter Ordnung.}
\end{center}

\section*{\S\,85. Ternäre Formen und algebraische Curven.}

Um im weiteren Verlauf der Darstellung einige von den mannigfaltigen und interessanten Anwendungen der Algebra auf geometrische Probleme vorführen zu können, ohne auf andere Hülfsmittel zu verweisen, werden hier die für diese Anwendungen nöthigen geometrischen Grundlagen kurz abgeleitet. Es wird dabei nur die Geometrie der Ebene betrachtet und von ihr nur das unmittelbar Erforderliche benutzt.\footnote{Ausführlicheres findet man in den Lehrbüchern der analytischen Geometrie, unter denen das Werk von George Salmon, \emph{Higher plane curves}, deutsch von Fiedler, hervorzuheben ist.}

Wir wenden homogene Dreieckscoordinaten an. Die Lage eines Punktes $x$ wird durch seine Abstände von den drei Seiten des Coordinatendreiecks, in drei beliebigen Maasseinheiten gemessen oder mit drei beliebigen constanten Factoren multiplicirt, bestimmt. Diese drei Grössen heissen die Coordinaten des Punktes $x$ und werden meist $x_1,x_2,x_3$ oder $y_1,y_2,y_3$ bezeichnet. Zwischen ihnen besteht zunächst eine lineare, nicht homogene Relation; mit Hülfe dieser Relation kann die Einheit linear und homogen durch $x_1,x_2,x_3$ dargestellt und jede nicht homogene Function in eine homogene verwandelt werden.

Die Coordinaten $x_1,x_2,x_3$ werden nunmehr als unabhängige Veränderliche betrachtet, wenn nur homogene Formen und Gleichungen benutzt werden. Nicht $x_1,x_2,x_3$ selbst, sondern nur die Verhältnisse $x_1:x_2:x_3$ bestimmen die Lage des Punktes. Die Werthcombination $x_1=x_2=x_3=0$ ist ausgeschlossen; die Werthe $hx_1,hx_2,hx_3$ bestimmen für jedes $h\ne0$ denselben Punkt.

Jede lineare Substitution
\begin{equation}
\begin{aligned}
 x_1&=\alpha^{(1)}_1y_1+\alpha^{(2)}_1y_2+\alpha^{(3)}_1y_3,\\
 x_2&=\alpha^{(1)}_2y_1+\alpha^{(2)}_2y_2+\alpha^{(3)}_2y_3,\\
 x_3&=\alpha^{(1)}_3y_1+\alpha^{(2)}_3y_2+\alpha^{(3)}_3y_3,
\end{aligned}
\tag{1}
\end{equation}
der deren Determinante
\begin{equation}
 r=\sum \pm \alpha^{(1)}_1\alpha^{(2)}_2\alpha^{(3)}_3
 \tag{2}
\end{equation}
von Null verschieden ist, kann doppelt gedeutet werden: als Abbildung zweier Punkte desselben Coordinatensystems oder als Coordinatentransformation eines und desselben Punktes. In den folgenden Betrachtungen wird die letztere Deutung festgehalten.

Irgend eine ternäre Form $n$ten Grades $f(x_1,x_2,x_3)$ stellt, gleich Null gesetzt, eine Curve $n$ter Ordnung dar, die mit einer geraden Linie $n$ Schnittpunkte hat. Durch (1) geht $f$ in eine Form desselben Grades in $y_1,y_2,y_3$ über,
\begin{equation}
 f(x_1,x_2,x_3)=\varphi(y_1,y_2,y_3),
 \tag{3}
\end{equation}
und $\varphi=0$ stellt dieselbe Curve in den neuen Coordinaten dar. Ableitung von (3) nach $y_1,y_2,y_3$ liefert
\begin{equation}
\begin{aligned}
 \varphi'(y_1)&=\alpha^{(1)}_1 f'(x_1)+\alpha^{(1)}_2 f'(x_2)+\alpha^{(1)}_3 f'(x_3),\\
 \varphi'(y_2)&=\alpha^{(2)}_1 f'(x_1)+\alpha^{(2)}_2 f'(x_2)+\alpha^{(2)}_3 f'(x_3),\\
 \varphi'(y_3)&=\alpha^{(3)}_1 f'(x_1)+\alpha^{(3)}_2 f'(x_2)+\alpha^{(3)}_3 f'(x_3).
\end{aligned}
\tag{4}
\end{equation}

\section*{\S\,86. Singuläre Punkte. Wendepunkte. Doppeltangenten.}

Wenn für einen Punkt $x$ die drei Gleichungen
\begin{equation}
 f'(x_1)=0,\qquad f'(x_2)=0,\qquad f'(x_3)=0
 \tag{1}
\end{equation}
gleichzeitig erfüllt sind, so liegt der Punkt nach dem Euler'schen Theorem
\begin{equation}
 nf=x_1f'(x_1)+x_2f'(x_2)+x_3f'(x_3)
 \tag{2}
\end{equation}
auf der Curve $f$. Die Gleichungen \S\,85, (4) zeigen zugleich, dass diese Eigenschaft durch lineare Transformationen erhalten bleibt. Nicht auf jeder Curve $n$ter Ordnung kommen solche Punkte vor. Durch Elimination von $x_1,x_2,x_3$ aus (1) erhält man eine Relation unter den Coefficienten von $f$; diese ganze rationale homogene Function der Coefficienten heisst die Discriminante von $f$ und ist nur bis auf einen numerischen Factor bestimmt. Die Punkte der Curve, deren Coordinaten (1) erfüllen, heissen singuläre Punkte.

Soll eine Curve unendlich viele singuläre Punkte besitzen, so müssen die partiellen Ableitungen einen gemeinsamen nicht constanten Theiler haben; aus denselben Erwägungen folgt dann, dass dieser Theiler auch mit $f$ einen gemeinsamen Theiler hat. Eine Curve mit unendlich vielen singulären Punkten ist also reducibel, und die Punkte liegen auf einem mehrfachen Bestandtheil der Curve.

Zur geometrischen Deutung ordnen wir $f$ nach Potenzen von $x_3$:
\begin{equation}
 f(x_1,x_2,x_3)=x_3^n f_0+x_3^{n-1}f_1+x_3^{n-2}f_2+\cdots,
 \tag{3}
\end{equation}
wobei $f_0,f_1,f_2,\ldots$ binäre Formen in $x_1,x_2$ sind und der Grad durch den Index angegeben wird. Liegt ein singulärer Punkt in der Ecke $I_3$, so müssen $f_0$ und $f_1$ verschwinden, und die Geraden durch $I_3$, welche die Curve berühren, ergeben sich aus $f_2=0$. Ist $f_2$ nicht identisch Null, so ist $I_3$ ein Doppelpunkt; ist auch $f_2=0$, so erhält man einen mehrfacheren singulären Punkt.

Ein Punkt einer Curve heisst Wendepunkt, wenn die Tangente die Curve dort dreipunktig berührt. Ist $I_3$ ein Wendepunkt und $x_1=0$ die Wendetangente, so hat $f$ die Gestalt
\begin{equation}
 f=x_1\varphi+x_3^2\psi,
 \tag{4}
\end{equation}
wobei $\varphi$ eine Form $(n-1)$ten und $\psi$ eine Form $(n-2)$ten Grades ist. Entsprechend erhält man für vierpunktige Berührung die stärkere Theilbarkeit der entsprechenden binären Form. Eine Gerade, welche die Curve in zwei getrennten Punkten berührt, heisst Doppeltangente; wählt man sie zu $x_3=0$ und die Berührungspunkte zu den auf $x_3=0$ liegenden Ecken, so muss in $f\vert_{x_3=0}$ sowohl $x_1^2$ als auch $x_2^2$ als Factor auftreten.

\section*{\S\,87. Covarianten der ternären Formen.}

Ist
\[
 f(x_1,x_2,x_3)=\sum a_{i_1i_2i_3}x_1^{i_1}x_2^{i_2}x_3^{i_3}
\]
eine ternäre Form, so heisst eine Form $C(x,a)$, die in den Variablen $x$ und in den Coefficienten $a$ homogen ist, eine Covariante von $f$, wenn sie bei der Transformation (1) der Variablen mit demselben Factor transformirt:
\begin{equation}
 C(y,b)=r^\lambda C(x,a).
 \tag{1}
\end{equation}
Ist die Form von den Variablen unabhängig, so ist sie eine Invariante. Der Exponent $\lambda$ ist das Gewicht.

Für alle ternären Formen vom Grade $>1$ tritt die Hesse'sche Determinante auf. Mit der hier gebrauchten Normirung schreiben wir
\begin{equation}
 H(x)=\det\bigl(f_{ij}\bigr)_{1\le i,j\le3},\qquad f_{ij}=f''(x_i,x_j),
 \tag{2}
\end{equation}
beziehungsweise, bis auf den im Buche gewählten numerischen Factor, $\Delta(x)$. Ausser dieser Covariante betrachtet Weber zwei weitere allgemeine Covarianten. Zunächst wird aus den zweiten Ableitungen von $f$ und den ersten Ableitungen der Hesse'schen Covariante eine Determinante gebildet:
\begin{equation}
 J(x)=\frac1{36}
 \left|\begin{array}{cccc}
 f_{11}&f_{12}&f_{13}&\Delta_1\\
 f_{21}&f_{22}&f_{23}&\Delta_2\\
 f_{31}&f_{32}&f_{33}&\Delta_3\\
 \Delta_1&\Delta_2&\Delta_3&0
 \end{array}\right|,
 \tag{3}
\end{equation}
wobei $\Delta_i=\Delta'(x_i)$. Dann folgt als dritte Covariante die Functional-Determinante der drei Formen $f,\Delta,J$:
\begin{equation}
 K(x)=\frac1{9}
 \left|\begin{array}{ccc}
 f_1&\Delta_1&J_1\\
 f_2&\Delta_2&J_2\\
 f_3&\Delta_3&J_3
 \end{array}\right|.
 \tag{4}
\end{equation}
Die Transformationsgesetze lauten in der hier gebrauchten Schreibweise
\begin{equation}
 \Delta(y)=r^2\Delta(x),\qquad J(y)=r^6J(x),\qquad K(y)=r^9K(x),
 \tag{5}
\end{equation}
und geben die betreffenden Gewichte an.

\section*{\S\,88. Die Hesse'sche Curve.}

Die durch $H=0$ dargestellte Curve heisst die Hesse'sche Curve der Curve $f=0$. Um ihren Zusammenhang mit den Wendepunkten zu sehen, sei im Punkte $I_3$ die Gleichung von $f$ in der Form
\begin{equation}
 f=x_3^n f_0+x_3^{n-1}f_1+x_3^{n-2}f_2+\cdots
 \tag{1}
\end{equation}
gegeben. Setzt man für einen Wendepunkt $I_3$ und eine Wendetangente $x_1=0$ die entsprechende Gestalt ein, so erhält man für die nach absteigenden Potenzen von $x_3$ geordnete Hesse'sche Covariante
\begin{equation}
 H=x_3^{3n-6}H_0+x_3^{3n-7}H_1+\cdots,
 \tag{2}
\end{equation}
mit
\begin{equation}
 H_0=-2(n-1)^2h^2c,
 \tag{3}
\end{equation}
und einem ausdrücklichen Ausdruck für $H_1$, der aus den ersten Gliedern der Entwickelung berechnet wird. Daraus folgt: die Hesse'sche Curve geht genau dann durch den Punkt $I_3$, wenn $c=0$, also genau dann, wenn $I_3$ ein Wendepunkt der Curve $f$ ist. Weiter liest man aus $H_1$ ab, wann die Wendetangente zugleich Tangente an die Hesse'sche Curve ist; bei einer Curve dritter Ordnung kann der dabei auftretende vierpunktige Fall nur eintreten, wenn die cubische Form reducibel wird.

Damit ist für eine nichtsinguläre cubische Curve die Bestimmung der Wendepunkte auf die Schnittpunkte von $f=0$ und $H=0$ zurückgeführt. Da beide Curven cubisch sind, erhält man neun Wendepunkte, sofern keine Besonderheiten eintreten.

\section*{\S\,89. Inflexionspunkte einer Curve dritter Ordnung.}

Es sei nun $f=0$ die Gleichung einer Curve dritter Ordnung ohne singulären Punkt. Man lege zwei Ecken des Coordinatendreiecks $I_1,I_2$ in zwei der neun Inflexionspunkte und wähle die zugehörigen Inflexionstangenten zu den Seiten $x_1=0$ und $x_2=0$. Dann muss sich $f$, wenn $x_1=0$ oder $x_2=0$ gesetzt wird, auf ein Glied $hx_3^3$ reduciren. Somit erhält $f$ die Gestalt
\begin{equation}
 f(x_1,x_2,x_3)=x_1x_2(a_1x_1+a_2x_2+a_3x_3)+hx_3^3.
 \tag{1}
\end{equation}
Die Gerade
\begin{equation}
 a_1x_1+a_2x_2+a_3x_3=0
 \tag{2}
\end{equation}
ist ebenfalls Inflexionstangente, und ihr Schnittpunkt mit $x_3=0$ ist der dritte Inflexionspunkt auf dieser Geraden.

Mit einer imaginären dritten Einheitswurzel $\varepsilon$ führt Weber zwei lineare Functionen $z_1,z_2$ ein, bestimmt durch
\begin{equation}
\begin{aligned}
 a_1x_1&=\varepsilon z_1+\varepsilon^2z_2-\tfrac13a_3x_3,\\
 a_2x_2&=\varepsilon^2z_1+\varepsilon z_2-\tfrac13a_3x_3,
\end{aligned}
\tag{3}
\end{equation}
und daraus folgt ferner
\begin{equation}
 -(a_1x_1+a_2x_2+a_3x_3)=z_1+z_2-\tfrac13a_3x_3.
 \tag{4}
\end{equation}
Durch Einsetzen in (1) ergibt sich nach geeigneter Multiplication mit constanten Factoren eine symmetrische Darstellung
\begin{equation}
 f=\varphi(y_1,y_2,y_3)=h(y_1^3+y_2^3+y_3^3)+6my_1y_2y_3.
 \tag{5}
\end{equation}
Damit folgt der Satz: Eine nichtsinguläre ternäre cubische Form lässt sich auf vier verschiedene Arten in die canonische Form (5) transformiren. Da auf jeder der Linien $y_1=0,y_2=0,y_3=0$ drei Inflexionspunkte liegen, ist das Problem der Inflexionspunkte wesentlich dasselbe wie die Transformation der cubischen Form auf die canonische Form.\footnote{Vgl. Newton und Maclaurin über die Wendepunkte der Curven dritter Ordnung, Hesse über die Elimination, und Aronhold über homogene Functionen dritten Grades.}

\section*{\S\,90. Transformation der cubischen Form auf die canonische Form.}

Für die canonische Form schreiben wir
\begin{equation}
 \varphi(y)=h\sum y_i^3+6my_1y_2y_3.
 \tag{1}
\end{equation}
Zur Bestimmung der Transformation werden die Covarianten der canonischen Form berechnet. Für die Hesse'sche Covariante ergibt sich
\begin{equation}
 \Delta(y)=-hm^2\sum y_i^3+(h^3+2m^3)y_1y_2y_3.
 \tag{2}
\end{equation}
Für die zweite Covariante erhält man in der Normirung
\begin{equation}
 J(x)=\frac1{36}
 \left|\begin{array}{cccc}
 f_{11}&f_{12}&f_{13}&\Delta_1\\
 f_{21}&f_{22}&f_{23}&\Delta_2\\
 f_{31}&f_{32}&f_{33}&\Delta_3\\
 \Delta_1&\Delta_2&\Delta_3&0
 \end{array}\right|,
 \tag{3}
\end{equation}
und für die canonische Form
\begin{align}
 J(y)=&-h^2(h^3+8m^3)^2\sum y_i^3y_j^3
      +3m^2(5h^6+26h^3m^3-4m^6)y_1^2y_2^2y_3^2\notag\\
 &+hm(2h^6+5h^3m^3+20m^6)y_1y_2y_3\sum y_i^3
      +9h^2m^6\left(\sum y_i^3\right)^2 .
 \tag{4}
\end{align}
Aus den Invariantenrelationen
\begin{equation}
 \varphi(y)=f,
 \qquad \Delta(y)=r^2\Delta,
 \qquad J(y)=r^6J
 \tag{5}
\end{equation}
folgen zuerst
\begin{align}
 (h^3+8m^3)y_1y_2y_3&=m^2f+r^2\Delta,
 \tag{6}\\
 h(h^3+8m^3)\sum y_i^3&=(h^3+2m^3)f-6mr^2\Delta.
 \tag{7}
\end{align}
Dann liefert (4)
\begin{equation}
\begin{aligned}
 h^2(h^3+8m^3)^2\sum y_i^3y_j^3={}&(2h^3+m^3)m^3f^2
 +m(2h^3-5m^3)r^2f\Delta\\
&+3m^2r^4\Delta^2-r^6J.
\end{aligned}
\tag{8}
\end{equation}

Die dritte Covariante wird in vereinfachter Form geschrieben als
\begin{equation}
 K(x)=\frac1{9}
 \left|\begin{array}{ccc}
 f_1&\Delta_1&J_1\\
 f_2&\Delta_2&J_2\\
 f_3&\Delta_3&J_3
 \end{array}\right|,
 \tag{9}
\end{equation}
und für die canonische Form findet man
\begin{equation}
 K(y)=-h^3(h^3+8m^3)^3
 \left|\begin{array}{ccc}
 1&y_1^3&y_1^6\\
 1&y_2^3&y_2^6\\
 1&y_3^3&y_3^6
 \end{array}\right|.
 \tag{10}
\end{equation}
Da $K(y)=r^9K$, folgt
\begin{equation}
 r^9K=h^3(h^3+8m^3)^3(y_1^3-y_2^3)(y_1^3-y_3^3)(y_2^3-y_3^3).
 \tag{11}
\end{equation}
Wenn kein singulärer Punkt vorhanden ist, kann $h^3+8m^3$ nicht verschwinden. Setzt man der Einfachheit halber $h=1$, so erhält man aus den Gleichungen (6), (7), (8) die cubische Gleichung
\begin{equation}
 u^3-Qu^2+Pu-R^3=0,
 \tag{12}
\end{equation}
deren Wurzeln $y_1^3,y_2^3,y_3^3$ sind, mit
\begin{align}
 P&={ (2+m^3)m^3f^2+(2-5m^3)mr^2f\Delta+3m^2r^4\Delta^2-r^6J\over (1+8m^3)^2},
 \tag{13}\\
 Q&={ (1+2m^3)f-6mr^2\Delta\over 1+8m^3},
 \tag{14}\\
 R&={m^2f+r^2\Delta\over 1+8m^3}.
 \tag{15}
\end{align}
Die Discriminante dieser cubischen Gleichung ist
\begin{equation}
 D_1={r^{18}K^2\over(1+8m^3)^6}.
 \tag{16}
\end{equation}

\section*{\S\,91. Invarianten der Curve dritter Ordnung. Realität.}

Zur Berechnung der Invarianten betrachtet Weber die Determinante der Hesse'schen Covariante der Form
\begin{equation}
 \Delta_{\lambda,\mu}=\lambda f+\mu\Delta.
 \tag{1}
\end{equation}
Sie lässt sich in der Form
\begin{equation}
 \Delta_{\lambda,\mu}(x)=Lf+M\Delta
 \tag{2}
\end{equation}
darstellen, wobei $L,M$ binäre Formen dritten Grades in $\lambda,\mu$ sind. Für die canonische Form erhält man, indem in \S\,90, (2) die Grössen
\[
 h,\\ m
\]
durch
\[
 h(\lambda-6m^2\mu),\qquad m\lambda+(h^3+2m^3)\mu
\]
ersetzt werden,
\begin{equation}
\begin{aligned}
\Delta_{\lambda,\mu}={}&-h(\lambda-6m^2\mu)[m\lambda+(h^3+2m^3)\mu]^2\sum y_i^3\\
&+\{h^3(\lambda-6m^2\mu)^3+2[m\lambda+(h^3+2m^3)\mu]^3\}y_1y_2y_3.
\end{aligned}
\tag{3}
\end{equation}
Daraus folgt
\begin{equation}
\begin{aligned}
 L&=-2\lambda^2\mu(h^3-m^3)m-\lambda\mu^2(h^6-20h^3m^3-8m^6)+8\mu^3m^2(h^3-m^3)^2,\\
 M&=\lambda^3+12\lambda^2\mu m(h^3-m^3)+2\mu^3(h^6-20h^3m^3-8m^6).
\end{aligned}
\tag{4}
\end{equation}
Man erhält also zwei Invarianten vierten und sechsten Grades, in canonischer Form
\begin{equation}
 m(h^3-m^3),
 \qquad h^6-20h^3m^3-8m^6.
 \tag{5}
\end{equation}
Bezeichnet man diese in der allgemeinen Form mit $S,T$ und setzt $h=1$, so ist
\begin{equation}
 m(1-m^3)=r^4S,
 \qquad 1-20m^3-8m^6=r^6T.
 \tag{6}
\end{equation}
Für $L,M$ erhält man dann
\begin{equation}
 L=-2S\lambda^2\mu-T\lambda\mu^2+8S^2\mu^3,
 \qquad M=\lambda^3+12S\lambda^2\mu+2T\mu^3.
 \tag{7}
\end{equation}
Aus (6) folgt für das Verhältnis $m^2:r^2$ die Gleichung
\begin{equation}
 27m^8+18Sm^4r^4+Tm^2r^6-S^2r^8=0.
 \tag{8}
\end{equation}
Durch die Substitution
\begin{equation}
 z={3m^2\over r^2}
 \tag{9}
\end{equation}
wird sie
\begin{equation}
 z^4+6Sz^2+Tz-3S^2=0.
 \tag{10}
\end{equation}
Die Discriminante dieser biquadratischen Gleichung ist
\begin{equation}
 D_2=-27(T^2+64S^3)^2=-27D^2,
 \tag{11}
\end{equation}
wobei $D=T^2+64S^3$. Für die canonische Form hat $D$ den Ausdruck
\begin{equation}
 h^3(h^3+8m^3)^3,
 \tag{12}
\end{equation}
und verschwindet daher nicht, solange die Curve keinen singulären Punkt hat. Aus (11) folgt, dass die biquadratische Gleichung bei reellen Curven eine negative Discriminante, also zwei reelle und zwei conjugirt imaginäre Wurzeln besitzt. Für die reelle positive Wurzel $z$ werden $m$ und $r$ reell, und daraus folgt der Satz:

\begin{quote}
Jede reelle ternäre cubische Form mit nicht verschwindender Discriminante kann durch eine reelle Transformation auf die reelle canonische Form gebracht werden.
\end{quote}

In den canonischen Coordinaten erhält man die neun Wendepunkte aus
\begin{equation}
 y_1^3+y_2^3+y_3^3+6my_1y_2y_3=0
 \tag{13}
\end{equation}
indem man einmal $y_1$, $y_2$, $y_3$ gleich Null setzt. Ist $\varepsilon$ eine imaginäre dritte Einheitswurzel, so ergeben sich die Coordinaten
\begin{equation}
\begin{array}{ccc}
(0,1,-1),&(0,1,-\varepsilon),&(0,1,-\varepsilon^2),\\
(-1,0,1),&(-\varepsilon,0,1),&(-\varepsilon^2,0,1),\\
(1,-1,0),&(1,-\varepsilon,0),&(1,-\varepsilon^2,0).
\end{array}
\tag{14}
\end{equation}
Je drei der in einer Zeile stehenden Punkte liegen auf einer der Geraden $y_1=0$, $y_2=0$, $y_3=0$. Die übrigen Tripel von Geraden entstehen in den drei folgenden Tabellen:
\begin{equation}
\begin{array}{ccc|ccc|ccc}
(0,1,-1)&(-1,0,1)&(1,-1,0)&(0,1,-1)&(-\varepsilon,0,1)&(1,-\varepsilon^2,0)&(0,1,-1)&(-\varepsilon^2,0,1)&(1,-\varepsilon,0)\\
(0,1,-\varepsilon)&(-\varepsilon,0,1)&(1,-\varepsilon,0)&(0,1,-\varepsilon)&(-\varepsilon^2,0,1)&(1,-1,0)&(0,1,-\varepsilon)&(-1,0,1)&(1,-\varepsilon^2,0)\\
(0,1,-\varepsilon^2)&(-\varepsilon^2,0,1)&(1,-\varepsilon^2,0)&(0,1,-\varepsilon^2)&(-1,0,1)&(1,-\varepsilon,0)&(0,1,-\varepsilon^2)&(-\varepsilon,0,1)&(1,-1,0)
\end{array}
\tag{15}
\end{equation}

Für die Realität ergeben sich, neben (14), die conjugirten Systeme
\begin{equation}
\begin{array}{ccc}
(0,1,-1),&(-\varepsilon,0,1),&(1,-\varepsilon^2,0)\\
(-1,0,1),&(1,-\varepsilon,0),&(0,1,-\varepsilon^2)\\
(1,-1,0),&(0,1,-\varepsilon),&(-\varepsilon^2,0,1),
\end{array}
\tag{16}
\end{equation}
und
\begin{equation}
\begin{array}{ccc}
(0,1,-1),&(-\varepsilon^2,0,1),&(1,-\varepsilon,0)\\
(-1,0,1),&(1,-\varepsilon^2,0),&(0,1,-\varepsilon)\\
(1,-1,0),&(0,1,-\varepsilon^2),&(-\varepsilon,0,1).
\end{array}
\tag{17}
\end{equation}
Eine übersichtlichere Bezeichnung erhält man, wenn $\xi$ und $\eta$ je ein volles Restsystem modulo $3$ durchlaufen. Die neun Wendepunkte werden dann durch die Paare $(\xi,\eta)$ bezeichnet, wobei $(\xi,\eta)$ und $(\eta,\xi)$ conjugirt imaginär sind und $(0,0),(1,1),(2,2)$ reell sind. Die vier Systeme von Geraden lauten
\begin{equation}
\begin{array}{ccc@{\qquad}ccc}
(0,0)&(1,2)&(2,1)&(0,0)&(1,1)&(2,2)\\
(1,1)&(2,0)&(0,2)&(1,2)&(2,0)&(0,1)\\
(2,2)&(0,1)&(1,0)&(2,1)&(0,2)&(1,0)\\[0.5em]
(0,0)&(2,0)&(1,0)&(0,0)&(0,2)&(0,1)\\
(1,1)&(0,1)&(2,1)&(1,1)&(1,0)&(1,2)\\
(2,2)&(1,2)&(0,2)&(2,2)&(2,1)&(2,0).
\end{array}
\tag{18}
\end{equation}
Daraus liest man, dass drei Wendepunkte $(\xi_1,\eta_1)$, $(\xi_2,\eta_2)$, $(\xi_3,\eta_3)$ genau dann auf einer Geraden liegen, wenn
\begin{equation}
 \xi_1+\xi_2+\xi_3\equiv0,
 \qquad
 \eta_1+\eta_2+\eta_3\equiv0
 \pmod 3.
 \tag{19}
\end{equation}
Die Bezeichnung der vier Geradensysteme kann auch aus dem Schema
\begin{equation}
\begin{array}{ccc}
(0,0)&(1,2)&(2,1)\\
(1,1)&(2,0)&(0,2)\\
(2,2)&(0,1)&(1,0)
\end{array}
\tag{20}
\end{equation}
gewonnen werden, indem man Zeilen, Columnen und die den positiven und negativen Gliedern der Determinantenbildung entsprechenden Combinationen bildet.


\end{document}
